\section{Revisão da Literatura}
\subsection{Envelhecimento cutâneo e estresse oxidativo}
O envelhecimento cutâneo resulta da ação conjunta de fatores intrínsecos, relacionados ao tempo e ao metabolismo celular, e de fatores extrínsecos, como radiação ultravioleta, poluição e tabagismo. A exposição acumulada pode aumentar a formação de espécies reativas de oxigênio (ERO); quando os sistemas antioxidantes não conseguem neutralizá-las, ocorre estresse oxidativo. Embora também participem da sinalização celular, em excesso as ERO danificam lipídios, proteínas e ácidos nucleicos e favorecem senescência e inflamação celular. Esses processos contribuem para rugas, perda de elasticidade e flacidez; a radiação ultravioleta também está associada à pigmentação irregular \parencite{liu2023,papaccio2022}.

A matriz extracelular (MEC) constitui uma rede de proteínas e outros componentes que sustenta e organiza os tecidos. Na derme, o colágeno confere resistência mecânica, a elastina permite a recuperação das fibras após a deformação e o ácido hialurônico retém água, contribuindo para a hidratação cutânea. As fibras de colágeno I e III suportam tensão, enquanto o colágeno IV integra a membrana basal que ancora a epiderme à derme. Com o envelhecimento, mudanças na síntese e na organização da MEC reduzem o suporte da pele e favorecem rugas, flacidez e perda de elasticidade \parencite{liu2025,papaccio2022}.

As metaloproteinases de matriz (MMPs) são proteases dependentes de zinco que participam da remodelação da MEC. A atividade dessas enzimas integra a renovação dos tecidos; quando a degradação supera a síntese, a matriz perde organização e resistência. A MMP-1 cliva os colágenos fibrilares I e III. As gelatinases MMP-2 e MMP-9 degradam gelatina e componentes da membrana basal, incluindo colágeno IV, na interface entre epiderme e derme. A elastase de neutrófilos humanos é uma serino-protease que degrada elastina, enquanto a hialuronidase humana 1 (HYAL-1) hidrolisa ácido hialurônico. Esses alvos, portanto, atuam sobre fibras de colágeno, fibras elásticas e glicosaminoglicanos da MEC \parencite{chao2007,liu2025,macdonald2002,papaccio2022}.

Além da remodelação da matriz, o estudo considera um alvo ligado à pigmentação. A melanina, produzida por melanócitos e distribuída aos queratinócitos, determina a coloração da pele. A proteína relacionada à tirosinase 1 (TYRP1) participa da melanogênese; a estrutura 5M8P contém um centro binuclear de zinco e foi determinada em complexo com tirosina \parencite{lai2017,papaccio2022}. O Quadro~\ref{quad:receptores_literatura} relaciona cada alvo à estrutura e às fontes bibliográficas utilizadas.

\begin{table}[H]
\centering
\small
\setlength{\tabcolsep}{3pt}
\renewcommand{\arraystretch}{1.15}
\quadrocaption[Receptores, códigos PDB e referências]{Receptores avaliados, códigos PDB e relações com processos cutâneos}
\label{quad:receptores_literatura}
\vspace{3pt}
\begin{tabular}{|C{2.8cm}|C{1.4cm}|C{6.3cm}|C{4.0cm}|}
\hline
\textbf{Receptor} & \textbf{PDB} & \textbf{Relação com a pele} & \textbf{Referências} \\
\hline
Colagenase humana (MMP-1) & 1HFC & Inicia a clivagem dos colágenos fibrilares I e III da derme. & \parencite{liu2025,spurlino1994} \\
\hline
Gelatinase A (MMP-2) & 8H78 & Degrada gelatina e colágeno IV da membrana basal. & \parencite{liu2025,takeuchi2023} \\
\hline
Gelatinase B (MMP-9) & 1GKC & Participa da degradação de gelatina e da remodelação da membrana basal. & \parencite{liu2025,rowsell2002} \\
\hline
Elastase de neutrófilos humanos & 1H1B & Degrada elastina das fibras elásticas. & \parencite{macdonald2002} \\
\hline
Hialuronidase humana 1 (HYAL-1) & 2PE4 & Hidrolisa o ácido hialurônico da matriz extracelular. & \parencite{chao2007,liu2025} \\
\hline
Proteína relacionada à tirosinase 1 (TYRP1) & 5M8P & Relaciona-se à melanogênese; apresenta centro binuclear de zinco. & \parencite{lai2017} \\
\hline
\end{tabular}
\end{table}

\subsection{Microalgas como fonte de compostos bioativos}
A variedade de moléculas produzidas por microalgas tem motivado pesquisas sobre aplicações biotecnológicas, inclusive em processos cutâneos ligados à matriz extracelular, à hidratação e à pigmentação. A interpretação dos resultados depende da espécie e do material analisado, pois biomassa, extratos, vesículas e substâncias poliméricas extracelulares (SPE) correspondem a materiais distintos \parencite{khemiri2023}.

Em \textit{T. chuii}, as pesquisas abrangem extratos, produtos da digestão simulada da biomassa e materiais lipídicos. Estudos avaliaram propriedades antioxidantes dos extratos e dos produtos da digestão, além da recuperação de ácidos graxos ômega-3 \parencite{cokdinleyen2024,moon2023,paterson2024}. Essa variedade de materiais contextualiza a investigação das SPE da espécie.

Hadiyanto \textit{et al.} (2025) obtiveram SPE de \textit{T. chuii} e avaliaram seu uso na agregação de microplásticos em água. O estudo documenta a produção de SPE pela espécie e uma aplicação ambiental desse material.

As vesículas extracelulares de \textit{T. chuii}, denominadas nanoalgossomos, são partículas delimitadas por membrana que transportam moléculas. Estudos avaliaram essas vesículas em modelos experimentais e em células cutâneas humanas \parencite{adamo2024,gargano2026}. Em conjunto, os trabalhos abordam materiais e finalidades distintos. A seção seguinte examina a composição das SPE e os componentes relatados nessa espécie e em materiais relacionados, que fundamentam a seleção dos compostos para prospecção.

Os componentes identificados em \textit{Tetraselmis} variam conforme a origem da amostra. A biomassa inclui as células e seus componentes; a teca é o revestimento que envolve a célula; já a SPE é liberada no meio de cultivo. Por isso, encontrar um açúcar na biomassa ou na teca não demonstra que ele também esteja presente na SPE secretada. O Quadro~\ref{quad:compostos_tetraselmis} apresenta separadamente esses materiais e os compostos relatados para cada espécie.

\begin{table}[H]
\centering
\small
\setlength{\tabcolsep}{3pt}
\renewcommand{\arraystretch}{1.12}
\quadrocaption[Compostos relatados em espécies de Tetraselmis]{Materiais e compostos relatados em espécies de \textit{Tetraselmis}, com distinção entre componentes celulares e extracelulares}
\label{quad:compostos_tetraselmis}
\vspace{3pt}
\begin{tabular}{|C{2.2cm}|C{3.2cm}|C{7.2cm}|C{2.7cm}|}
\hline
\textbf{Espécie} & \textbf{Material analisado} & \textbf{Componentes ou achados relatados} & \textbf{Referência} \\
\hline
\textit{T. chuii} & Biomassa & Glicose, galactose, manose, arabinose e ramnose foram descritas na composição de açúcares da biomassa. & Brown (1991, \textit{apud} Pereira \textit{et al.}, 2019) \\
\hline
\textit{T. chuii} & SPE recuperada do cultivo & A produção de SPE foi quantificada; a amostra seca apresentou carboidratos. O estudo não detalhou a composição monossacarídica dessa SPE. & Hadiyanto \textit{et al.} (2025) \\
\hline
\textit{T. suecica} & Exopolissacarídeos extraídos do cultivo & Foram identificados glicose, ácido glucurônico, manose, galactose, arabinose, ribose, ramnose, fucose e xilose; as proporções variaram com a condição de cultivo e o tipo de exopolissacarídeo extraído. & \textcite{parra2020} \\
\hline
\textit{T. striata} & Parede celular (teca) & Foram identificados KDO, KDO metilado e ácido 3-deoxi-lixo-2-heptulosárico (Dha). & \textcite{becker1989} \\
\hline
\textit{T. tetrahele} & Extrato de biomassa & Foram identificados compostos fenólicos, incluindo ácido gálico e ácido ferúlico. & \textcite{mai2026} \\
\hline
\end{tabular}
\end{table}


\subsection{Substâncias poliméricas extracelulares e seu potencial}
As SPE são materiais liberados pelas microalgas, presentes ao redor das células ou recuperados do meio de cultivo. Sua fração polissacarídica contém unidades monossacarídicas ligadas entre si. A espécie e as condições de cultivo influenciam sua composição, enquanto o método de extração determina quais frações são recuperadas \parencite{delattre2016,he2025}.

Grupos hidroxila, carboxila e sulfato influenciam as propriedades dos polissacarídeos. A sequência dos açúcares, as ramificações e as substituições também definem a estrutura da cadeia \parencite{delattre2016,he2025}.

Os estudos publicados abrangem materiais e espécies distintos, da biomassa de \textit{T. chuii} a polissacarídeos extracelulares e vesículas de outras microalgas. Essa diversidade define o contexto para interpretar a composição e as atividades relatadas nos trabalhos relacionados ao tema.

Estudos com polissacarídeos sulfatados de \textit{Chlamydomonas reinhardtii} relataram atividade antioxidante e inibição de enzimas associadas à matriz cutânea \parencite{falcao2020}. Em \textit{Porphyridium} sp., foram caracterizadas unidades de glicose sulfatadas, que oferecem contexto para interpretar derivados sulfatados incluídos na triagem \parencite{geresh2009}. Essas evidências dizem respeito a outras espécies e frações.

Os 17 compostos selecionados para prospecção incluem açúcares, derivados de açúcares e fenólicos descritos em diferentes espécies e materiais. O Quadro~\ref{quad:ligantes_spe_literatura} relaciona cada composto às evidências publicadas em \textit{Tetraselmis} e em outras microalgas.

\begin{table}[H]
\centering
\footnotesize
\setlength{\tabcolsep}{3pt}
\renewcommand{\arraystretch}{1.08}
\quadrocaption[Compostos selecionados para prospecção e evidências em Tetraselmis]{Compostos selecionados para prospecção e ocorrência relatada em espécies de \textit{Tetraselmis}}
\label{quad:ligantes_spe_literatura}
\vspace{3pt}
\begin{tabular}{|C{3.0cm}|C{8.2cm}|C{4.0cm}|}
\hline
\textbf{Ligante selecionado} & \textbf{Material, espécie e evidência relatada} & \textbf{Referência} \\
\hline
Ácido ferúlico & Extrato de biomassa de \textit{T. tetrahele}. & \textcite{mai2026} \\
\hline
Ácido galacturônico & Exopolissacarídeos de \textit{T. suecica}; também relatado na teca de \textit{T. striata}. & \textcite{parra2020,becker1989} \\
\hline
Ácido gálico & Extrato de biomassa de \textit{T. tetrahele}. & \textcite{mai2026} \\
\hline
Ácido glucurônico & Exopolissacarídeos de \textit{T. suecica}. & \textcite{parra2020} \\
\hline
D-Galactose & Biomassa de \textit{T. chuii}; exopolissacarídeos de \textit{T. suecica}; teca de \textit{T. striata}. & Brown (1991, \textit{apud} Pereira \textit{et al.}, 2019); \textcite{parra2020,becker1989} \\
\hline
D-Glicose & Biomassa de \textit{T. chuii}; exopolissacarídeos de \textit{T. suecica}; teca de \textit{T. striata}. & Brown (1991, \textit{apud} Pereira \textit{et al.}, 2019); \textcite{parra2020,becker1989} \\
\hline
D-Gulose & Teca de \textit{T. striata}. & \textcite{becker1989} \\
\hline
D-Manose & Biomassa de \textit{T. chuii}; exopolissacarídeos de \textit{T. suecica}. & Brown (1991, \textit{apud} Pereira \textit{et al.}, 2019); \textcite{parra2020} \\
\hline
D-Xilose & Exopolissacarídeos de \textit{T. suecica}. & \textcite{parra2020} \\
\hline
Dha & Teca de \textit{T. striata}. & \textcite{becker1989} \\
\hline
Galactose-6-sulfato & Não localizada em \textit{Tetraselmis}. $\beta$-D-galactose-6-sulfato foi identificada nos funoranos de algas vermelhas do gênero \textit{Gloiopeltis}; $\alpha$-L-galactose-6-sulfato também foi descrita em \textit{Acanthophora spicifera}. & \textcite{tuvikene2015funorans,duarte2004} \\
\hline
Glicose-6-sulfato & Em polissacarídeo sulfatado de \textit{T. suecica}, o sulfato foi proposto no O-6 de resíduos de glicose; o estudo não relata o monossacarídeo isolado. & \textcite{li2026tetraselmis} \\
\hline
KDO & Teca de \textit{T. striata}. & \textcite{becker1989} \\
\hline
L-Arabinose & Biomassa de \textit{T. chuii}; exopolissacarídeos de \textit{T. suecica}; teca de \textit{T. striata}. & Brown (1991, \textit{apud} Pereira \textit{et al.}, 2019); \textcite{parra2020,becker1989} \\
\hline
L-Fucose & Exopolissacarídeos de \textit{T. suecica}. & \textcite{parra2020} \\
\hline
L-Ramnose & Biomassa de \textit{T. chuii}; exopolissacarídeos de \textit{T. suecica}. & Brown (1991, \textit{apud} Pereira \textit{et al.}, 2019); \textcite{parra2020} \\
\hline
$\alpha$-Kdo & Na teca de \textit{T. striata}, Kdo integra os polissacarídeos da parede; em bactérias Gram-negativas, resíduos $\alpha$-Kdo ligam o lipídio A ao núcleo do lipopolissacarídeo (LPS). Não foi localizado relato de $\alpha$-Kdo isolado em \textit{Tetraselmis}. & \textcite{becker1989,becker1998tetraselmis,ovchinnikova2016} \\
\hline
\end{tabular}
\end{table}
\Needspace{12\baselineskip}
\subsection{Ligantes comerciais e coocristalizados selecionados}
Os ligantes comerciais foram incluídos para comparar as poses e os contatos previstos para os compostos selecionados para prospecção. O Quadro~\ref{quad:ligantes_comerciais_literatura} reúne os 17 ligantes comerciais selecionados, indica quais também são coocristalizados e apresenta as fontes das estruturas. A L-tirosina foi incluída como ligante coocristalizado suplementar para TYRP1, fora do conjunto de 34 ligantes da triagem.

\begin{table}[H]
\centering
\footnotesize
\setlength{\tabcolsep}{3pt}
\renewcommand{\arraystretch}{1.15}
\quadrocaption[Ligantes comerciais e coocristalizados selecionados]{Ligantes comerciais e coocristalizados considerados no estudo e respectivas referências}
\label{quad:ligantes_comerciais_literatura}
\vspace{3pt}
\begin{tabular}{|C{3.0cm}|C{3.3cm}|C{4.0cm}|C{4.0cm}|}
\hline
\textbf{Ligante} & \textbf{Categoria} & \textbf{Coocristalizado em} & \textbf{Referência da estrutura} \\
\hline
ARP-100 & Comercial & -- & \parencite{pubchem2023} \\
\hline
Ácido ascórbico & Comercial & -- & \parencite{pubchem2023} \\
\hline
Ácido kójico & Comercial & -- & \parencite{pubchem2023} \\
\hline
Apigenina & Comercial & -- & \parencite{pubchem2023} \\
\hline
Arbutina & Comercial & -- & \parencite{pubchem2023} \\
\hline
Batimastat & Comercial & -- & \parencite{pubchem2023} \\
\hline
Cromoglicato & Comercial & -- & \parencite{pubchem2023} \\
\hline
Doxiciclina & Comercial & -- & \parencite{pubchem2023} \\
\hline
GW475151 & Coocristalizado & Elastase (1H1B) & \parencite{macdonald2002} \\
\hline
Hidroquinona & Comercial & -- & \parencite{pubchem2023} \\
\hline
L2U & Coocristalizado & MMP-2 (8H78) & \parencite{takeuchi2023} \\
\hline
Marimastat & Comercial & -- & \parencite{pubchem2023} \\
\hline
NFH & Coocristalizado & MMP-9 (1GKC) & \parencite{rowsell2002} \\
\hline
PLH & Coocristalizado & MMP-1 (1HFC) & \parencite{spurlino1994} \\
\hline
Quercetina & Comercial & -- & \parencite{pubchem2023} \\
\hline
Sivelestat & Comercial & -- & \parencite{pubchem2023} \\
\hline
Tropolona & Comercial & -- & \parencite{pubchem2023} \\
\hline
L-tirosina & Coocristalizado & TYRP1 (5M8P) & \parencite{lai2017} \\
\hline
\end{tabular}
\end{table}

As estruturas químicas dos ligantes comerciais sem complexo cristalográfico foram consultadas no PubChem. Os artigos indicados para os coocristalizados documentam as estruturas dos complexos e sua posição cristalográfica; essa informação permite comparar as poses, sem atribuir a todos os ligantes atividade experimental nos seis receptores.

\subsection{Docking molecular como ferramenta de triagem virtual}
No \textit{docking} molecular, o algoritmo posiciona o ligante no sítio de ligação do receptor e explora diferentes orientações e conformações. Cada posição resultante é chamada de pose. Uma função de pontuação atribui um valor a cada pose para estimar como o ligante se ajusta ao sítio. Os resultados variam conforme o procedimento de busca, que define quais posições e conformações são exploradas, e conforme a representação molecular adotada, que descreve a estrutura tridimensional e as propriedades químicas atribuídas ao ligante e ao receptor \parencite{eberhardt2021vina,trott2010}.

O AutoDock Vina é o programa de busca; nele, podem ser usadas diferentes funções de pontuação, incluindo Vina e AD4. A função Vina combina termos que representam contato estérico, interações hidrofóbicas e ligações de hidrogênio, além de considerar a flexibilidade do ligante. AD4 inclui termos de eletrostática e dessolvatação. Para alvos que contêm zinco, AD4Zn ajusta os parâmetros usados para calcular as interações com esse metal \parencite{eberhardt2021vina,morris2009autodock,santos2014zn,trott2010}. Cada receptor é analisado com a função de pontuação empregada em seu cálculo.

Nos resultados, FEB é a pontuação reportada em kcal/mol, denominada em inglês \textit{free energy of binding} (energia livre de ligação). A função de pontuação usa a FEB para ordenar as poses: valores mais negativos indicam poses mais bem avaliadas no cálculo. Assim, em cada receptor, a FEB organiza os ligantes conforme a pontuação obtida com a função e o protocolo adotados \parencite{warren2006}.

Os contatos previstos ajudam a interpretar as poses. O PLIP identifica tipos de contato entre proteína e ligante com base em critérios geométricos, e o PyMOL permite visualizar os complexos e examinar a posição do ligante no sítio \parencite{salentin2015plip,pymol}.

O desvio quadrático médio (RMSD, do inglês \textit{root-mean-square deviation}) compara as coordenadas de átomos correspondentes. Nesta análise, a conformação do ligante na pose obtida pelo \textit{docking} é comparada à do ligante coocristalizado. Os dois ligantes são alinhados entre si, e o RMSD resume a diferença entre as posições dos átomos correspondentes após esse ajuste. Valores menores indicam maior semelhança entre as conformações \parencite{cole2005}.

A análise conjunta da FEB, dos contatos e do RMSD permite priorizar os compostos para futuras avaliações experimentais de atividade, segurança e permeação cutânea.
